\section{Mise en place des plots et des cultures}
\subsection{Mise en place des plots}
La mise en place des plots renvoie à la délimitation et au nivellement des 6 plots après un labour profond. Pour délimiter nos plots nous avons utilisé la méthode 3-4-5. C'est est une méthode simple basée sur le théorème de Pythagore pour tracer des angles droits (90 degrés). En agriculture, elle est très utile pour délimiter correctement les planches, les parcelles et les lignes de culture. Dans notre cas, pour l'obtention de tous les sommets, nous avons utilisé un piquet pour marquer un point A, puis avec le ruban-mètre nous avons mesuré 3~m pour avoir le point B. À partir de A, dans une direction perpendiculaire, nous mesurons 4~m pour y placer le piquet C, puis à partir de ces 3 points nous projetons avec des arcs de cercle afin d'obtenir les 4 points déterminant ainsi la surface totale de nos plots expérimentaux.
Après la délimitation des plots, nous avons élevé les ados, puis nous avons bien nivelé et enfin nous avons fixé les ados en les arrosant. Ce schéma ci-dessous est une parfaite illustration de la méthode 3-4-5.

\begin{figure}[H]
    \centering
    \imgstd{images/16_methode_3_4_5.png}
    \caption{Illustration de la méthode 3-4-5 pour la délimitation de la surface des plots. }
    \label{fig:plot3-4-5}
\end{figure}

La figure ci-dessous présente la configuration obtenue après utilisation de la méthode 3-4-5 pour la délimitation de notre surface.

\begin{figure}[H]
    \centering
    \imgstd{images/17_fin_confection.jpg}
    \caption{Fin de confection des 6 plots. }
    \label{fig:confection3}
\end{figure}

\subsection{Mise en place des cultures sur les plots}

\begin{itemize}[label=$\bullet$]
   	 \item \textbf{Organisation des cultures sur les plots}

     Nous avons 6 plots constitués de 3 plots témoins et de 3 plots tests. Par annotation, de la gauche vers la droite, nous avons sur la première ligne les plots de gombo, au milieu les plots de navet et à droite les plots de laitue. Chaque plot test est aligné sur la même ligne avec son témoin de la même spéculation. Le plot gombo test est nommé GOM TS et son témoin GOM T, le plot test navet est nommé NAV TS et son témoin NAV T, et enfin, le plot laitue test est nommé LAI TS et son témoin LAI T.
	\end{itemize}
\begin{itemize}[label=$\bullet$]
   	 \item \textbf{Mise en place des cultures}

     Nous avons réalisé un semis direct de gombo sur les plots GOM TS et GOM T, et de navet sur les plots NAV TS et NAV T, le soir du 25 juin 2026. Pour les plots de laitue, un repiquage a été effectué le soir du 25 juillet 2026.
	\end{itemize}

\begin{itemize}[label=$\bullet$]
   	 \item \textbf{Méthodes de semis et de repiquage}

     Pour le gombo, nous avons semé par poquets. Dans chaque poquet, nous avons mis une ou deux graines, avec une profondeur égale à deux fois la taille de la graine. Le semis du navet est de type semis en ligne, avec une ou deux graines par trou et une distanciation de 15 cm. La laitue a également été repiquée en ligne, avec espacement de 15 cm.
	\end{itemize}





